% 图 F-07 Q1 求解原理示意图 — 紧凑版
\documentclass[border=5pt]{article}
\usepackage[margin=0pt,paperwidth=25cm,paperheight=10cm]{geometry}
\pagestyle{empty}
\usepackage{fontspec}
\usepackage{amsmath}
\usepackage{ctex}
\newcommand{\fontdir}{../../../00_知识库/论文写作/LaTeX模板_2026国赛版/fonts/}
\setCJKfamilyfont{song}[
  Path = \fontdir, UprightFont = SourceHanSerifCN-Regular.otf,
  BoldFont = SourceHanSerifCN-Bold.otf, AutoFakeBold = true,
]{SourceHanSerifCN}
\newcommand*{\song}{\CJKfamily{song}}
\usepackage{tikz}
\usetikzlibrary{arrows.meta,positioning,calc}

\definecolor{primary}{HTML}{1F3A93}
\definecolor{accent}{HTML}{D35400}
\definecolor{green}{HTML}{27AE60}
\definecolor{gray}{HTML}{95A5A6}
\definecolor{textDark}{HTML}{333333}
\definecolor{gridC}{HTML}{CCCCCC}
\definecolor{beige}{HTML}{FFF8EE}
\definecolor{ctrlFill}{HTML}{E8F5E9}

\tikzset{arrH/.style={-{Stealth[length=2.2mm,width=2.0mm]},line width=1.2pt,draw=accent}}
\tikzset{arrM/.style={-{Stealth[length=2.2mm,width=2.0mm]},line width=1.2pt,draw=primary}}
\tikzset{arrG/.style={-{Stealth[length=2.2mm,width=2.0mm]},line width=1.0pt,draw=accent}}

\begin{document}\song
\begin{tikzpicture}[scale=0.85]

% ====== (a) 圆柱截面与边界 ======
\draw[fill=primary!4,draw=primary,line width=1.3pt,rounded corners=2pt]
  (-5.5,-3.2) rectangle (3.5,4.8);
\node[font=\song\small\bfseries,text=textDark] at (-1,4.3)
  {(a) 圆柱截面与边界条件};

% 横向椭圆（再下移）
\draw[fill=beige,draw=accent,line width=1.2pt] (-1,0.1) ellipse (2.5 and 0.9);

% 中心点 r=0
\draw[fill=textDark] (-1,0.1) circle (1.8pt);
\node[font=\song\footnotesize,text=textDark] at (-1.4,-0.4) {$r=0$};

% 径向箭头
\draw[draw=textDark,->,line width=0.7pt] (-1,0.1) -- (1.5,0.1);
\node[font=\song\footnotesize,text=textDark] at (0.2,0.35) {$r$};
\node[font=\song\footnotesize,text=accent] at (2.2,0.6) {$R_0=2$ cm};

% Robin 条件
\node[font=\song\footnotesize,text=textDark,align=center] at (-1,3.5)
  {$-k\partial_r T = h(T_R-T_\infty),\ -D\partial_r C = k_m(C_R-C_\infty)$};

% 热流 / 湿流箭头
\draw[arrH] (-1.8,1.0) -- (-1.8,2.2);
\node[font=\song\footnotesize,text=accent] at (-1.8,2.45) {热流};
\draw[arrM] (-0.3,1.0) -- (-0.3,2.2);
\node[font=\song\footnotesize,text=primary] at (-0.3,2.45) {湿流};

% 对称条件
\node[font=\song\footnotesize,text=textDark,align=center] at (-1,-1.3)
  {$r=0$：$\partial_r T = \partial_r C = 0$};

% 网格
\foreach \x in {0,0.6,...,2.4}
  \draw[gridC,line width=0.3pt] (-1,0.1) circle (\x cm and \x*0.36 cm);

% ====== (b) 径向控制体 ======
\draw[fill=primary!4,draw=primary,line width=1.3pt,rounded corners=2pt]
  (4.5,-3.2) rectangle (13.5,4.8);
\node[font=\song\small\bfseries,text=textDark] at (9,4.3)
  {(b) 径向控制体 $(N=80,\ \Delta r=0.25\ \mathrm{mm})$};

% 控制体方块（整体下移，拉开与标题的距离）
\fill[beige,draw=accent,line width=1.0pt] (5.5,0.4) rectangle (7.0,2.9);
\node[font=\song\footnotesize,text=accent] at (6.25,3.2) {中心半格};
\fill[ctrlFill,draw=green,line width=0.8pt] (7.0,0.4) rectangle (8.5,2.9);
\fill[ctrlFill,draw=green,line width=0.8pt] (8.5,0.4) rectangle (10.0,2.9);
\fill[ctrlFill,draw=green,line width=0.8pt] (10.0,0.4) rectangle (11.5,2.9);
\fill[ctrlFill,draw=green,line width=1.0pt] (11.5,0.4) rectangle (12.3,2.9);
\node[font=\song\footnotesize,text=green!50!black] at (11.9,3.2) {表面半格};

% 分隔线
\foreach \xi in {7.0,8.5,10.0,11.5}
  \draw[draw=textDark,line width=0.5pt] (\xi,-0.1) -- (\xi,2.9);

\node[font=\song\footnotesize,text=textDark,align=center] at (9,-0.5)
  {$r_{i-1/2}$（内界面）\quad $r_{i+1/2}$（外界面）};
\node[font=\song\footnotesize,text=textDark,align=center] at (9,-1.2)
  {输出点 = 计算点（$0.1\ \mathrm{cm}=4\Delta r$）};

% ====== (c) 时空离散与推进 ======
\draw[fill=white,draw=gray,line width=0.9pt,rounded corners=2pt]
  (14.5,-3.2) rectangle (23.0,4.8);
\node[font=\song\small\bfseries,text=textDark] at (18.75,4.3)
  {(c) 时空离散与推进};

% 坐标轴
\draw[->,line width=0.8pt] (15.5,-2.5) -- (15.5,3.8) node[above,font=\song\footnotesize] {$t$};
\draw[->,line width=0.8pt] (15.5,-2.5) -- (22.5,-2.5) node[right,font=\song\footnotesize] {$r$};

% 网格
\foreach \x in {16.2,16.9,...,22.0}
  \draw[gridC,line width=0.3pt] (\x,-2.2) -- (\x,3.2);
\foreach \y in {-2.2,-1.4,...,3.2}
  \draw[gridC,line width=0.3pt] (15.5,\y) -- (21.5,\y);

% 推进箭头
\draw[arrG] (16.2,-1.8) -- (16.2,0.2);
\node[font=\song\footnotesize,text=accent,rotate=90] at (15.7,-0.8) {推进};

% 步长文字
\node[font=\song\small,text=textDark,align=left] at (19.2,2.8)
  {输出步: $\Delta t_o = 1\ \mathrm{s}$\\[2pt]
   热步子步: $\Delta t_T = 1/32\ \mathrm{s}$\\[2pt]
   湿步子步: $\Delta t_C = 0.1\ \mathrm{s}$};

% 时间层级标注
\node[font=\song\footnotesize,text=gray] at (20.5,-1.5) {$K$ 时层};
\node[font=\song\footnotesize,text=textDark] at (20.5,1.0) {$K{+}1$ 时层};
\draw[gridC,line width=0.5pt,dashed] (15.5,-0.2) -- (21.5,-0.2);

\end{tikzpicture}
\end{document}